\section{Validation on a Pretrained Language Model}\label{sec:realdata}
% Every number below is a macro from generated/numbers_lm.tex (scripts/make_lm_tables.py).
% Write the interpretation sentences ONLY after the Colab run; state the result whichever way it goes.
A controlled Boolean setting establishes correctness but may make attribution too easy. We therefore apply the
identical pipeline (matched-budget PLT and CLT, feature-to-feature graphs with error nodes, the same depth
metric) to two-hop factual recall in \LMModel{}~\cite{lindsey2025biology}:
\emph{``Fact: the capital of the state containing Dallas is''} $\to$ \emph{Austin}. The intermediate entity
(\emph{Texas}) is known, and a one-hop control prompt with the same answer
(\emph{``Fact: the capital of Texas is''}) removes exactly one inferential step. We keep prompts whose answer
ranks in the model's top ten and whose state capital the model also knows, leaving \LMNTwoHop{} two-hop and
\LMNOneHop{} one-hop prompts. Both transcoders use the same features per layer, TopK $k$, steps, data and seed;
held-out normalised MSE is \LMNMSEPLT{} (PLT) and \LMNMSECLT{} (CLT).

\paragraph{V1: is the depth metric sensitive to depth on real data?}
If $\Dfeat$ measures inferential depth, it should be larger from the city tokens of a two-hop prompt than from
the state tokens of its one-hop control. Paired by state, $\Dfeat(\text{2-hop})-\Dfeat(\text{1-hop})$ is
\LMVOnePLT{} [\LMVOneLoPLT{}, \LMVOneHiPLT{}] for the PLT ($p=\LMVOnePPLT$) and
\LMVOneCLT{} [\LMVOneLoCLT{}, \LMVOneHiCLT{}] for the CLT ($p=\LMVOnePCLT$).

\paragraph{V2: are the state features causally used?}
State-selective features are chosen leave-one-out, so the prompt being tested never influences the selection.
Zeroing them lowers the answer logit by \LMVTwoPLT{} [\LMVTwoLoPLT{}, \LMVTwoHiPLT{}] (PLT) and
\LMVTwoCLT{} [\LMVTwoLoCLT{}, \LMVTwoHiCLT{}] (CLT) \emph{more} than zeroing the same number of random active
features in the same layers ($p=\LMVTwoPPLT$, $p=\LMVTwoPCLT$).

\paragraph{V3: does the graph's claim match the intervention?}
Across prompts, the graph-claimed mediation through state features (IMF) correlates with the measured ablation
drop with Spearman $\rho=\LMVThreePLT$ (PLT, $p=\LMVThreePPLT$) and $\rho=\LMVThreeCLT$ (CLT, $p=\LMVThreePCLT$).

\paragraph{V4: CLT versus PLT.}
On the same two-hop prompts, CLT minus PLT $\Dfeat$ is \LMDfeatDiffCLTminusPLT{} ($p=\LMDfeatDiffP$);
Table~\ref{tab:lm} reports skip mass, IMF, ablation drop and error influence.

% AUTO-GENERATED by scripts/make_lm_tables.py
\begin{table}[t]\centering\scriptsize
\caption{Real-data validation on two-hop factual recall. Mean [95\% bootstrap CI]; sign-flip $p$.}\label{tab:lm}
\begin{tabular}{llp{3.6cm}ll}\toprule
Test & Kind/hop & Quantity & Mean [CI] & $p$ (n)\\\midrule
V1 depth sensitivity & clt & D feat(2hop) - D feat(1hop), per state & 0.09 [-0.06, 0.22] & 0.344 (n=6)\\
V1 depth sensitivity & clt & abs-normalised D feat(2hop) - D feat(1hop), per state & -0.02 [-0.04, 0.00] & 0.219 (n=6)\\
V2 causal specificity & clt & logit drop: state features - random features & 4.33 [2.27, 6.52] & 0.016 (n=7)\\
V3 graph vs intervention & clt & Spearman(IMF, ablation drop) & 0.64 [--, --] & 0.119 (n=7)\\
V1 depth sensitivity & plt & D feat(2hop) - D feat(1hop), per state & 0.28 [0.17, 0.38] & 0.031 (n=6)\\
V1 depth sensitivity & plt & abs-normalised D feat(2hop) - D feat(1hop), per state & -0.13 [-0.20, -0.07] & 0.031 (n=6)\\
V2 causal specificity & plt & logit drop: state features - random features & 0.83 [-0.85, 2.93] & 0.594 (n=7)\\
V3 graph vs intervention & plt & Spearman(IMF, ablation drop) & 0.57 [--, --] & 0.180 (n=7)\\
V4 clt minus plt & 2hop & root signed D feat & -0.16 [-0.35, -0.02] & 0.059 (n=10)\\
V4 clt minus plt & 1hop & root signed D feat & -0.03 [-0.15, 0.09] & 0.608 (n=12)\\
V4 clt minus plt & 2hop & skip mass & 0.07 [0.07, 0.08] & 0.002 (n=10)\\
V4 clt minus plt & 1hop & skip mass & 0.06 [0.05, 0.06] & 0.000 (n=12)\\
V4 clt minus plt & 2hop & IMF h1 & 0.04 [0.01, 0.08] & 0.047 (n=10)\\
V4 clt minus plt & 2hop & state ablation logit drop & 0.72 [-0.08, 1.64] & 0.188 (n=10)\\
V4 clt minus plt & 2hop & error influence & -0.08 [-0.09, -0.07] & 0.002 (n=10)\\
V4 clt minus plt & 1hop & error influence & -0.07 [-0.08, -0.06] & 0.000 (n=12)\\
\bottomrule\end{tabular}\end{table}


\paragraph{Scope.} The real-model transcoders are trained by us at a matched budget but are far smaller than
production transcoders; error nodes carry \LMErrInflPLT{} (PLT) and \LMErrInflCLT{} (CLT) of the logit's
influence, and the replacement model without error terms keeps the original top-1 token on a fraction
\LMTopKeptPLT{} (PLT) and \LMTopKeptCLT{} (CLT) of prompts. InterpBench~\cite{gupta2024interpbench}
(semi-synthetic SIIT transformers with known circuits) is the natural next step between our certified
Boolean models and a pretrained LM; the evaluation code is included but was not run for this submission.
